\section{Discussion}
\label{sec:discussion}

\subsection{RQ1: Residual Topology and Performance}
\label{sec:rq1}

The main benchmark argues against a universal residual topology: MatrixRes is
the most frequent winner with the best mean rank, but all five model families
win at least two of the 24 dataset--operator combinations, and the GCNConv slice
also has dataset-dependent winners. The useful residual neighborhood therefore
depends on dataset-specific tolerance for branch specialization, residual
traffic, and optimization complexity. This is consistent with the broader
literature, where changing the propagation operator, normalization, or pooling
design shifts the preferred
architecture~\cite{xu2019gin,hamilton2017inductive,ying2018hierarchical}. The
present study adds residual topology itself as another source of that shift.
Matrix-style reuse can be read as a local residual stencil over the
branch-by-layer grid, analogous in spirit to structured neighborhood
mixing~\cite{abu2019mixhop} but applied to hidden-state reuse rather than graph
adjacency; its advantage is architectural organization, not a new
message-passing primitive.

\subsection{RQ2: Branch Count and Mechanisms}
\label{sec:rq2}

Increasing \(B\) adds branches, residual paths, and optimization interactions at
once, so it is not a simple capacity knob. PROTEINS tolerates a moderate budget
(peaks at \(B=4\)--\(6\)) whereas DD prefers \(B=1\)--\(2\). The mechanism
metrics are consistent with this split: on PROTEINS the best settings combine
increased branch diversity with reduced branch cosine, while on DD the best
settings have nearly saturated cosine and CKA, so additional residual traffic is
unlikely to create new functions. Diversity may rise with \(B\) but does not by
itself improve classification when branch similarity stays high or gradients
weaken. This rise-then-fall behaviour is a practical warning for residual GNN
design, since residual connections are often introduced to make deeper or wider
models easier to optimize~\cite{he2016deep}. The balance between
functional diversity and residual interference differs sharply between two
datasets that are both routinely treated as protein-structure benchmarks.

\subsection{RQ3: When Matrix-Style Reuse Helps}
\label{sec:rq3}

Sparse/gated reuse is not uniformly better than dense matrix reuse: MatrixRes
wins more combinations overall. MatrixResGated remains useful as a controlled
way to suppress weak residual coefficients, and the mechanism table confirms
that it keeps fewer active pathways than dense MatrixRes while the gate stays at
a moderate value. Its value is not that sparsity always improves accuracy, but
that filtering helps when matrix reuse would otherwise inject too much
low-value traffic. This places it between two familiar ideas: like residual
gated graph networks it learns how strongly historical information should be
injected~\cite{bresson2017residual}, and like sparsification methods it reduces
the influence of weak pathways~\cite{rong2019dropedge}. The implementation is
deliberately simple---a scalar gate and soft-shrink filtering rather than a
large routing network---which keeps the ablation interpretable. Candidate
settings must still be validated across folds: the DD hidden-dimension candidate
survives that check and yields a smaller, more efficient model, whereas the
PROTEINS sparsity candidate does not and should not be claimed as a gain.

The synthetic benchmark separates class granularity from real-dataset
idiosyncrasies. Simple residual choices remain competitive at two and three
classes, which guards against the overbroad claim that matrix-style reuse is
always needed. From four classes onward the pattern reverses: MatrixResGated has
the best average and normalized accuracy, MatrixRes has the best macro-F1 and
the smallest degradation from \(C=2\) to \(C=8\). Distribution-level results
refine this: MatrixRes leads in the degree- and density-driven BA, ER, and
regular-degree families, whose labels separate related structural regimes rather
than unrelated categories, while community-driven SBM is less favourable. Taken
together, matrix-style reuse is conditionally effective rather than universal,
and is most useful when the task requires separating several related structural
regimes under finer discrimination.

\subsection{Limitations}
\label{sec:limitations}

Several factors bound these findings. The benchmark uses a fixed graph-level
pooling and classification head, so alternative readouts may interact with
residual topology. The compared operators are canonical message-passing
backbones rather than recent graph transformers or local-global
hybrids~\cite{dwivedi2020generalization}, which is a deliberate control choice
but means the study is not a state-of-the-art comparison. The sensitivity scan is fold-0 only, and the candidate reruns cover
MatrixResGated on PROTEINS and DD alone. ENZYMES has low absolute accuracy and
high variability. The mechanism analysis is correlational: branch diversity,
cosine similarity, CKA, and gradient norms support a coherent explanation but do
not prove causality. Finally, the benchmark intentionally avoids task-specific
chemical or biological feature engineering, so absolute accuracies should not be
read as optimized state-of-the-art molecular prediction results, and the
synthetic families stress-test residual topology rather than substitute for
broader real multi-class benchmarks. Future work should extend the comparison to
larger datasets, richer readouts, recent local-global backbones, and learned
residual-neighborhood policies that replace the hand-specified vertical,
horizontal, and matrix stencils.
